\documentclass{article}

\usepackage[a4paper,margin=1in]{geometry}
\usepackage{amsmath,amssymb}
\usepackage{booktabs}

\title{\textbf{CSE400 – Fundamentals of Probability in Computing}}
\author{\textbf{Kensi Patel - AU2440102}}
\date{Lecture 19}

\begin{document}

\maketitle

\textbf{Lecture 19: N–Variate Joint PMF / PDF / CDF}

Instructor: Dhaval Patel, PhD

\section*{1. Joint Distribution of $N$ Random Variables}

\textbf{Definition}
\begin{itemize}
    \item For random variable $X_1,X_2,\ldots,X_N$, the vector representation is
\[\mathbf{X}=\begin{bmatrix}X_1\\X_2\\\vdots\\X_N\end{bmatrix}\]

    \item The joint distribution describes the probability behaviour of all variables together.
    \item It can be represented using:
    \begin{itemize}
        \item Joint PMF (discrete case)
        \item Joint PDF (continuous case)
        \item Joint CDF
    \end{itemize}
\end{itemize}
(Definition slide — page 17)

\section*{2. Joint PMF}

For discrete random variables:
\[P_{X_1,X_2,\ldots,X_N}(x_1,x_2,\ldots,x_N)=P(X_1=x_1,x_2,\ldots,X_N=x_N)\]
(Page 22)

\section*{3. Joint CDF}

\[F_{X_1,\ldots,X_N}(x_1,\ldots,x_N)=P(X_1\le x_1,\ldots,X_N\le x_N)\]
(Page 23)

\section*{4. Joint PDF}

Joint PDF is obtained by differentiating joint CDF:

\[f_{X_1,\ldots,X_N}(x_1,\ldots,x_N)=\frac{\partial^N}{\partial x_1\cdots \partial x_N}F_{X_1,\ldots,X_N}(x_1,\ldots,x_N)\]
(Page 24)

\section*{5. Independence of Random Variables}

Let
\[\mathbf{X}=[X_1,X_2,\ldots,X_N]^T\]

Variables are independent if:
\[f_{\mathbf{X}}(x_1,\ldots,x_N)=\prod_{i=1}^{N} f_{X_i}(x_i)\]
(Page 25)

\section*{Example 1 — Jointly Gaussian Random Variables}

Given:
\[\mu_X=2,\quad \sigma_X^2=1\]
\[\mu_Y=-3,\quad \sigma_Y^2=4\]
\[\text{cov}(X,Y)=-1\]

Find:
\begin{itemize}
    \item[(a)] Joint PDF
    \item[(b)] Marginal PDFs
    \item[(c)] $P(X<0)$
\end{itemize}
(Page 4)

\subsection*{(a) Construction of Joint Gaussian PDF}

\subsection*{Step 1: Correlation Coefficient}

\[\rho_{XY}=\frac{\text{cov}(X,Y)}{\sigma_X\sigma_Y}\]
From variances: 
\[\sigma_X=\sqrt{1}= 1,\qquad \sigma_Y=\sqrt{4}= 2\]
Substitute:
\[\rho_{XY}=\frac{-1}{(1)(2)}= -\frac{1}{2}\]
(Pages 7–8)

\subsection*{Step 2: General Joint Gaussian Formula}

\[f_{XY}(x,y)=\frac{1}{2\pi\sigma_X\sigma_Y\sqrt{1-\rho^2}}\exp\left(-\frac{1}{2(1-\rho^2)}\left[\left(\frac{x-\mu_X}{\sigma_X}\right)^2-2\rho\frac{(x-\mu_X)(y-\mu_Y)}{\sigma_X\sigma_Y}+\left(\frac{y-\mu_Y}{\sigma_Y}\right)^2\right]\right)\]
(Page 5)

\subsection*{Step 3: Substitute Values}

\[\mu_X=2,\ \sigma_X=1,\ \mu_Y=-3,\ \sigma_Y=2,\ \rho=-0.5\]
Final Joint PDF:
\[f_{XY}(x,y)=\frac{1}{2\pi\sqrt{3}}\exp\left(-\frac{2}{3}\left[(x-2)^2+(x-2)(y+3)+\frac{(y+3)^2}{4}\right]\right)\]
(Page 10)

\subsection*{(b) Marginal PDFs}

Marginals are obtained by integration:

\[f_X(x)=\int_{-\infty}^{\infty} f_{XY}(x,y)\,dy\]
\[f_Y(y)=\int_{-\infty}^{\infty} f_{XY}(x,y)\,dx\]
(Page 11)

\subsection*{Key Property}

If variables are jointly Gaussian:
\[X\sim N(2,1),\qquad Y\sim N(-3,4)\]
Thus:
\[f_X(x)=\frac{1}{\sqrt{2\pi}}e^{-\frac{(x-2)^2}{2}}\]
\[f_Y(y)=\frac{1}{2\sqrt{2\pi}}e^{-\frac{(y+3)^2}{8}}\]
(Page 12)

\subsection*{(c) Probability $P(X<0)$}

From marginal:
\[X\sim N(2,1)\]

\subsection*{Step 1: CDF Relation}
\[P(X<0)=F_X(0)\]
(Page 13)

\subsection*{Step 2: Standardization}
\[Z=\frac{X-\mu_X}{\sigma_X}\]
\[P(X<0)=P\!\left(Z<\frac{0-2}{1}\right)=P(Z<-2)=\Phi(-2)\]
(Page 14)

\subsection*{Step 3: Symmetry of Normal CDF}
\[\Phi(-x)=1-\Phi(x)=P(Z>x)\]
Hence,
\[P(Z<-2)=P(Z>2)\]
(Page 15)

\subsection*{Step 4: Relation to Q-Function}
\[Q(x)=P(Z>x)\]

\textbf{Final Answer}

\[P(X<0)=Q(2)\]
(Page 16)

\section*{Application — Large MIMO System}
\begin{itemize}
    \item System has $M$ transmit antennas and $N$ receive antennas.
    \item Each channel path coefficient $h_{ij}$ is modeled as a random variable.
    \item Entire system represented by channel matrix:
\end{itemize}
\[H=\begin{bmatrix}
h_{11} & \cdots & h_{1N} \\
\vdots & \ddots & \vdots \\
h_{M1} & \cdots & h_{MN}
\end{bmatrix}\]
(Pages 18–21)

\section*{Example 2 — Uniformly Distributed 3D Vector}

Region:
\[|x_1|\le1,\quad |x_2|\le1,\quad |x_3|\le1\]

Joint PDF:
\[f(x_1,x_2,x_3)=
\begin{cases}
C,& \text{inside cube}\\
0,& \text{otherwise}
\end{cases}\]
(Page 26)

\subsection*{(a) Find Constant}
Total probability = 1:
\[\int_{-1}^{1}\int_{-1}^{1}\int_{-1}^{1}C\,dx_1dx_2dx_3=1\]
\[C(2)(2)(2)=1\]
\[C=\frac{1}{8}\]
(Page 28)

\subsection*{(b) Marginal Joint PDF $f_{X_1,X_2}$}

\[f_{X_1,X_2}(x_1,x_2)=\int_{-1}^{1}\frac{1}{8}dx_3=\frac{1}{4}\]
(Page 29)

\subsection*{(c) Marginal PDF $f_{X_1}$}

\[f_{X_1}(x_1)=\int_{-1}^{1}\int_{-1}^{1}\frac{1}{8}dx_2dx_3=\frac{1}{2}\]
(Page 31)

\subsection*{(d) Conditional PDFs}

\[f_{X_1|X_2,X_3}=\frac{1/8}{1/4}=\frac{1}{2}\]
(Page 33)

\subsection*{(e) Conditional PDFs}
\[f_{X_1,X_2|X_3}=\frac{1/8}{1/2}=\frac{1}{4}\]
(Page 35)

\subsection*{(f) Independence Check}
LHS:
\[f_{X_1,X_2,X_3}=\frac{1}{8}\]
RHS:
\[f_{X_1}f_{X_2}f_{X_3}=\frac{1}{2}\cdot\frac{1}{2}\cdot\frac{1}{2}=\frac{1}{8}\]
Since LHS = RHS → variables are independent.

(Page 39)


\section*{Example 3 — Trivariate Exponential-type Distribution}

Joint PDF:
\[f_{XYZ}(x,y,z)=
\begin{cases}
Ke^{-(ax+by+cz)},& x,y,z>0\\
0,& \text{otherwise}
\end{cases}\]
(Page 40)

\subsection*{(a) Find $K$}

\[\int_{0}^{\infty}\int_{0}^{\infty}\int_{0}^{\infty}K e^{-(ax+by+cz)}\,dx\,dy\,dz = 1\]

\textbf{Factorization:}

\[K\left(\frac{1}{a}\right)\left(\frac{1}{b}\right)\left(\frac{1}{c}\right)=1\]
\[K=abc\]
(Page 42)

\subsection*{(b) Marginal Joint PDF}

\[f_{XY}(x,y)=\int_{0}^{\infty} abc\,e^{-(ax+by+cz)}\,dz\]
\[=ab\,e^{-(ax+by)}\]
(Page 44)

\subsection*{(c) Marginal PDF}

\[f_X(x)=\int_{0}^{\infty} ab\,e^{-(ax+by)}\,dy\]
\[= a e^{-ax}\]
(Page 46)

\subsection*{(d) Independence Check}

\textbf{LHS:}
\[f_{XYZ}=abc\,e^{-(ax+by+cz)}\]
\textbf{RHS:}
\[f_X f_Y f_Z=(a e^{-ax})(b e^{-by})(c e^{-cz})\]
Thus:
\[\text{LHS}=\text{RHS}\]
Hence $X,Y,Z$ are independent.

(Page 50)

\end{document}